\originalsection{黏性、边界条件与耗散}
把黏性应力单独记为$\tau_{ij}$, 则Cauchy应力为$-p\delta_{ij}+\tau_{ij}$, 动量通量为$\rho v_iv_j+p\delta_{ij}-\tau_{ij}$. 对各向同性牛顿流体, 令$S_{ij}=(\partial_iv_j+\partial_jv_i)/2$, $d=\nabla\cdot\vect v$, 有$\tau_{ij}=2\eta(S_{ij}-d\delta_{ij}/3)+\zeta d\delta_{ij}$. 对称速度梯度而非反对称部分进入此式, 是因为纯匀速平移或刚体转动没有邻层形变, 不应产生这种摩擦. 系数$\eta,\zeta\ge0$使
\[
\tau_{ij}\partial_jv_i
=2\eta\sum_{ij}(S_{ij}-d\delta_{ij}/3)^2+\zeta d^2\ge0.
\]
这表示机械能变成热的速率. 对不可压缩常密度流体, 无体力时, 动能密度$k=\rho v^2/2$满足
\[
\partial_tk+\partial_j\bigl[(k+p)v_j-v_i\tau_{ij}\bigr]
=-\tau_{ij}\partial_jv_i.
\]
积分后左侧的边界项是机械能通量. 若边界速度为零或在无穷远足够快衰减, 总动能导数为$-2\eta\int S_{ij}S_{ij}\dd V\le0$. 开放流动不能删掉边界通量, 例如压力泵持续输入能量才维持管流. 无热传导时$\rho T Ds/Dt=\tau_{ij}\partial_jv_i$, 故理想流体的熵守恒已被黏性产生的热改变.

无体力的常密度不可压缩流还可从Navier--Stokes取散度, 得压强Poisson方程$\nabla^2p=-\rho(\partial_iv_j)(\partial_jv_i)$. 这里用了$\partial_iv_i=0$, 说明压强是满足不可压缩约束所需的场, 不能先任意给压强再独立给速度. 固体壁面在黏性模型中通常要求无滑动$\vect v=\vect v_{\mathrm{wall}}$, 同时固定法向与切向速度. 这是与理想流体仅要求法向不可穿透的实质差别.

\originalsection{板间流与相似流动}
两平板在$y=0,h$, 定常充分发展速度为$v_x(y)$, 沿$x$压强梯度为$k<0$. 对流项$v_x\partial_xv_x$为零, Navier--Stokes给$\eta v_x''=k$. 两壁静止时$v_x(0)=v_x(h)=0$, 积分得$v_x=-ky(h-y)/(2\eta)$, 单位宽度流量为$-kh^3/(12\eta)$. 若$k=0$而上板以$U$运动, 条件为$v_x(0)=0$, $v_x(h)=U$, 得Couette流$v_x=Uy/h$. 一般二者可在这个线性约化模型中叠加. 对底壁的流向剪应力大小为$-kh/2$（静止两板）或$\eta U/h$（纯Couette流）; 必须明确是流体对壁还是壁对流体, 两者方向相反.

当几何只有一个尺度$L$、入流速度$U$且其他无量纲比也相同, 定常不可压缩流可写$\vect v/U=\vect h(\vect r/L,\mathrm{Re})$, $(p-p_\infty)/(\rho U^2)=g(\vect r/L,\mathrm{Re})$. 同Re的两个装置因而有相同无量纲场, 称动力相似. 这并不使Re成为任何几何下唯一参数: 若存在两个长度、可压缩性、自由表面或粗糙度, 其比值或其他无量纲数也须相同. 小Re问题的自然压力尺度改为$\eta U/L$, 它与惯性压力尺度相差$1/\mathrm{Re}$.

\originalsection{球体Stokes流的边值计算}
现取半径$a$固定球体, 无穷远$\vect v=\vect u=U\hat{\vect z}$, 球面无滑动$\vect v=0$. 常密度定常Navier--Stokes中惯性项尺度为$U^2/a$, 黏性项为$\nu U/a^2$, 比值为$\mathrm{Re}=Ua/\nu$. 当它很小, 主导边值问题为
\[
\nabla p=\eta\nabla^2\vect v,\qquad \nabla\cdot\vect v=0,
\qquad \vect v|_{r=a}=0,\quad\vect v\to\vect u\ (r\to\infty).
\]
散度给$\nabla^2p=0$. 边界问题对$\vect u$线性, 且只有球对称几何及入流方向, 取压强的衰减偶极形式$p=p_\infty+C\eta a(\vect u\cdot\vect r)/r^3$. 同样把速度写成
\[
\vect v=f(r)\vect u+k(r)(\vect u\cdot\vect r)\vect r.
\]
这两个矢量结构分别沿入流方向和径向; $f$无量纲, $k$量纲为长度的负二次方. 对该式求散度得$f'+r^2k'+4rk=0$. 再求Laplacian, 比较两个结构, 得
\[
\begin{aligned}
k''+\frac{6k'}r&=-\frac{3Ca}{r^5},\\
f''+\frac{2f'}r+2k&=\frac{Ca}{r^3}.
\end{aligned}
\]
第一式的解为$k=Ca/(2r^3)+A+B/r^5$. 远场要求$r^2k\to0$, 故$A=0$; 球面在所有角度为零要求$k(a)=f(a)=0$, 故$B=-Ca^3/2$. 将所得$k$代入散度式并以$f(\infty)=1$积分, 得
\[
k=\frac{Ca}{2r^3}\left(1-\frac{a^2}{r^2}\right),\qquad
f=1+\frac{Ca}{2r}+\frac{Ca^3}{6r^3}.
\]
最后$f(a)=0$给$C=-3/2$. 代回另一个ODE也成立, 故没有遗漏动量方程. 速度和压力为
\[
\begin{aligned}
v_r&=U\left(1-\frac{3a}{2r}+\frac{a^3}{2r^3}\right)\cos\theta,\\
v_\theta&=-U\left(1-\frac{3a}{4r}-\frac{a^3}{4r^3}\right)\sin\theta,\\
p&=p_\infty-\frac{3\eta Ua}{2r^2}\cos\theta.
\end{aligned}
\]
两速度分量在球面都为零, 远处恢复均匀流. 涡量为$\Omega_\phi=-3Ua\sin\theta/(2r^2)$, 并非理想位势流. 原稿第137页在$f$的积分常数和末次边界条件的文字上有排印矛盾, 上面保留$f(\infty)=1$, $f(a)=0$独立核对了全部条件.

\originalsection{压力与剪应力共同产生阻力}
球面外法向为$\hat{\vect r}$, 流体对球的牵引为$\vect t=(-pI+\tau)\hat{\vect r}$. 在球面,
\[
\tau_{rr}=2\eta\partial_rv_r=0,\qquad
\tau_{r\theta}=\eta\left(\frac1r\partial_\theta v_r+
\partial_rv_\theta-\frac{v_\theta}r\right)
=-\frac{3\eta U}{2a}\sin\theta.
\]
由于$\hat{\vect r}\cdot\hat{\vect z}=\cos\theta$, $\hat{\vect\theta}\cdot\hat{\vect z}=-\sin\theta$, 球受流向力为
\[
F_z=a^2\int_0^{2\pi}\!\int_0^\pi
[-p\cos\theta+\tau_{rr}\cos\theta-\tau_{r\theta}\sin\theta]
\sin\theta\dd\theta\dd\phi.
\]
均匀压强$p_\infty$贡献为零. 余下压力项为$2\pi\eta aU$, 切向黏性项为$4\pi\eta aU$, 总和为$6\pi\eta aU$. 因此静止流体中以$\vect V$运动的球受到
\[
\vect F_{\mathrm{drag}}=-6\pi\eta a\vect V.
\]
不能只积分压强而漏掉剪应力, 也不能只计算剪应力而忽略黏性流的非均匀压强. 稳态小Re解给阻力的主导项; 非定常附加质量、历史力、壁面效应与有限Re修正没有包含在该式中.

\begin{center}
\begin{tikzpicture}\begin{axis}[width=.46\textwidth,height=5cm,axis equal image,xmin=-4,xmax=4,ymin=-2.6,ymax=2.6,xtick={-4,-2,0,2,4},ytick={-2,0,2},xlabel={$z/a$},ylabel={$s/a$},title={理想位势流},tick label style={font=\scriptsize},title style={font=\small}]
\addplot[blue!65!black,no marks,thin] table[x=x,y=y,col sep=comma]{figures/sphere-ideal-0.csv};
\addplot[blue!65!black,no marks,thin] table[x=x,y expr=-\thisrow{y},col sep=comma]{figures/sphere-ideal-0.csv};
\addplot[blue!65!black,no marks,thin] table[x=x,y=y,col sep=comma]{figures/sphere-ideal-1.csv};
\addplot[blue!65!black,no marks,thin] table[x=x,y expr=-\thisrow{y},col sep=comma]{figures/sphere-ideal-1.csv};
\addplot[blue!65!black,no marks,thin] table[x=x,y=y,col sep=comma]{figures/sphere-ideal-2.csv};
\addplot[blue!65!black,no marks,thin] table[x=x,y expr=-\thisrow{y},col sep=comma]{figures/sphere-ideal-2.csv};
\addplot[blue!65!black,no marks,thin] table[x=x,y=y,col sep=comma]{figures/sphere-ideal-3.csv};
\addplot[blue!65!black,no marks,thin] table[x=x,y expr=-\thisrow{y},col sep=comma]{figures/sphere-ideal-3.csv};
\addplot[black,fill=gray!20,domain=0:360,samples=73] ({cos(x)},{sin(x)});
\draw[->,thick] (axis cs:-3.7,2.0)--(axis cs:-2.7,2.0) node[midway,above] {$\vect u$};
\end{axis}\end{tikzpicture}\hfill
\begin{tikzpicture}\begin{axis}[width=.46\textwidth,height=5cm,axis equal image,xmin=-4,xmax=4,ymin=-2.6,ymax=2.6,xtick={-4,-2,0,2,4},ytick={-2,0,2},xlabel={$z/a$},ylabel={$s/a$},title={Stokes黏性流},tick label style={font=\scriptsize},title style={font=\small}]
\addplot[blue!65!black,no marks,thin] table[x=x,y=y,col sep=comma]{figures/sphere-stokes-0.csv};
\addplot[blue!65!black,no marks,thin] table[x=x,y expr=-\thisrow{y},col sep=comma]{figures/sphere-stokes-0.csv};
\addplot[blue!65!black,no marks,thin] table[x=x,y=y,col sep=comma]{figures/sphere-stokes-1.csv};
\addplot[blue!65!black,no marks,thin] table[x=x,y expr=-\thisrow{y},col sep=comma]{figures/sphere-stokes-1.csv};
\addplot[blue!65!black,no marks,thin] table[x=x,y=y,col sep=comma]{figures/sphere-stokes-2.csv};
\addplot[blue!65!black,no marks,thin] table[x=x,y expr=-\thisrow{y},col sep=comma]{figures/sphere-stokes-2.csv};
\addplot[blue!65!black,no marks,thin] table[x=x,y=y,col sep=comma]{figures/sphere-stokes-3.csv};
\addplot[blue!65!black,no marks,thin] table[x=x,y expr=-\thisrow{y},col sep=comma]{figures/sphere-stokes-3.csv};
\addplot[black,fill=gray!20,domain=0:360,samples=73] ({cos(x)},{sin(x)});
\draw[->,thick] (axis cs:-3.7,2.0)--(axis cs:-2.7,2.0) node[midway,above] {$\vect u$};
\end{axis}\end{tikzpicture}
\end{center}

图的流线由轴对称流函数等值线实际求根得到. 理想球体流$\psi=Ur^2\sin^2\theta(1-a^3/r^3)/2$; Stokes流$\psi=Ur^2\sin^2\theta(1-3a/(2r)+a^3/(2r^3))/2$. 二者都不穿过球面, 但只有后者切向无滑动. 此图帮助读出绕流几何; 阻力结论来自上面的应力积分, 不能由流线形状直接猜得.

\originalsection{随Re改变的绕流与湍流概述}
原课最后讨论增大Re后的现象, 没有把完整非线性绕流解写成公式. 当$\mathrm{Re}\ll1$, 上述流线前后对称; Re增到约1时惯性修正使尾部开始不同; 数十到数百的范围可出现分离涡与回流, 进而失去定常性并周期脱落; 更大Re时尾流含许多尺度并呈不规则变化. 原稿按球半径的特征尺度举$10$--$100$、$150$--$1000$、$10^3$--$2\times10^5$等近似流态范围, 这些是所讨论装置的现象尺度, 不是由Re定义推出的普适临界数. 球面粗糙度、扰动、实验条件及用半径还是直径定义Re都会影响转变.

高Re时边界层转为湍流可能推迟分离, 使尾流变窄并降低总阻力; 因而“黏性耗散增强”与“某物体阻力总是增大”不是同一命题. 原课以球面凹坑和球类表面缝线为例说明这种作用. 管内湍流的瞬时速度不规则, 平均速度常比层流抛物线更平坦; 不能继续把充分发展层流公式当作其完整描述.

三维强湍流常有跨尺度能量传递: 大尺度输入的机械能传到更小涡, 最终由黏性变成热. 这是一种统计图景, 并不声称每个涡都按同一时序重复. 原课第7章进一步用量纲分析介绍惯性区间与Kolmogorov谱, 将在非线性章继续. 本册不以这些概述冒充湍流闭合理论或一般Navier--Stokes正则性证明.
